\chapter{Análisis del problema}\label{cap:analisis}

\section{Introducción}

En este capítulo se analiza el problema que motiva el desarrollo del sistema \textbf{Sevillaneando}, así como los requisitos que debe cumplir la solución propuesta.

Actualmente, la información sobre eventos culturales y de ocio en la ciudad de Sevilla se encuentra dispersa en múltiples plataformas, muchas de ellas centradas en eventos de gran escala. Esto dificulta el descubrimiento de actividades emergentes o menos conocidas, reduciendo su visibilidad y limitando la participación de los usuarios. Personalmente, me cuesta mucho encontrar actividades culturales si no es porque algún conocido sepa de ellas, por lo que pensé que sería una buena idea tener un sitio donde poder encontrarlas.

También, las soluciones existentes suelen carecer de funcionalidades de personalización, lo que impide ofrecer una experiencia adaptada a los intereses y preferencias de cada usuario. Además, no suelen ofrecer un servicio social donde cada uno pueda aportar algo, por lo que es difícil dar a conocer nuevos lugares y eventos, y sobretodo, tener a alguien con quien comunicarte para obtener más información y para conocer a alguien que frecuente esos sitios.

Ante esta situación, se plantea el desarrollo de una aplicación móvil que centralice la información de eventos, permita la participación activa y social de los usuarios y ofrezca recomendaciones personalizadas basadas en su comportamiento, ubicación e intereses.

\section{Requisitos de información}

Los requisitos de información describen los datos que el sistema debe gestionar para poder ofrecer sus funcionalidades.

En el sistema \textbf{Sevillaneando}, los principales elementos de información son los siguientes:

\begin{itemize}
    \item \textbf{Usuarios}: información básica del usuario, incluyendo nombre de usuario, correo electrónico, contraseña, imagen de perfil, ubicación e intereses. La preferencia de tema visual se gestiona en el lado del cliente y no se persiste en la base de datos. También se almacena el rol del usuario: usuario, administrador o moderador.

    \item \textbf{Roles y permisos}: información asociada al tipo de usuario y a las acciones que puede realizar dentro del sistema, como crear eventos (rol de usuario), moderarlos (rol de moderador), gestionar usuarios o administrar categorías (rol de administrador).

    \item \textbf{Categorías}: clasificación de los eventos según su temática. Estas categorías son gestionadas por el administrador, es decir, este puede crearlas, editarlas y eliminarlas. Luego, son utilizadas por los usuarios para filtrar eventos y definir intereses para tenerlas en cuenta para las recomendaciones.

    \item \textbf{Eventos}: cada evento contiene título, descripción, categoría, ubicación, fecha, precio, imágenes, estado (aprobado, rechazado o pendiente), creador y visibilidad pública o privada.

    \item \textbf{Eventos privados}: información relativa al acceso restringido, incluyendo códigos o enlaces de acceso o invitación y usuarios autorizados, además de todos los datos que contiene une evento.

    \item \textbf{Moderación de eventos}: datos necesarios para controlar el estado de revisión de los eventos públicos propuestos por usuarios, incluyendo aprobación, rechazo y posibles modificaciones realizadas por el moderador de edición de eventos.

    \item \textbf{Asistencia a eventos}: relación entre usuarios y eventos que permite conocer qué usuarios están apuntados a cada evento y obtener un listado de eventos posterior.

    \item \textbf{Eventos guardados}: relación entre usuarios y eventos almacenados como guardados para su consulta en un listado posterior.

    \item \textbf{Calendario}: información derivada de los eventos a los que un usuario se ha apuntado y que tienen una fecha definida.

    \item \textbf{Mensajes}: datos asociados a la comunicación entre usuarios, tanto en chats de eventos como en conversaciones privadas, que puede ser fecha, emisor, destinatario (si el mensaje es privado), evento (si el mensaje es en el chat de evento), contenido e imagen (si se envía un mensaje de este tipo).

    \item \textbf{Imágenes}: fotografías asociadas a eventos, chats o perfiles de usuario, incluyendo imágenes subidas en tiempo real durante los eventos.

    \item \textbf{Valoraciones y reseñas}: puntuaciones y comentarios que los usuarios realizan sobre eventos y rutas.

    \item \textbf{Notificaciones}: avisos enviados al usuario, como aprobación o rechazo de eventos públicos creados/editados por el usuario, recordatorios de eventos apuntados en fechas cercanas y eventos cercanos según su ubicación.

    \item \textbf{Recomendaciones}: información utilizada para sugerir eventos en función de intereses, ubicación, eventos vistos, eventos guardados, eventos apuntados, eventos compartidos, fecha y valoraciones tanto propias como globales.

    \item \textbf{Rutas}: agrupaciones de eventos generadas automáticamente o creadas por los usuarios, con información sobre una secuencia de eventos , el orden entre ellos, franjas horarias, distancia entre eventos. Además incluye valoraciones en el caso de las rutas creadas como públicas y puntuación en base a los gustos mencionados en las recomendaciones en las rutas recomendadas personalmente.

    \item \textbf{Filtros de búsqueda}: información necesaria para buscar eventos por estado, categoría, título, descripción, precio, cercanía y radio de distancia.

    \item \textbf{Datos geográficos}: coordenadas y ubicación de usuarios y eventos, utilizadas para mostrar eventos en el mapa, calcular la cercanía, calcular las recomendaciones y generar las rutas recomendadas.
\end{itemize}

\vspace{0.1cm}

Estos datos se almacenan en una base de datos relacional, concretamente \textbf{PostgreSQL con PostGIS}. Esta permite gestionar de forma eficiente las relaciones entre entidades, así como realizar consultas tanto convencionales como geoespaciales. En particular, las capacidades de PostGIS se utilizan para operaciones como el cálculo de distancias, filtrado por radio, búsqueda de eventos cercanos, visualización en mapas y generación de rutas, entre otras funcionalidades basadas en localización.

\section{Requisitos funcionales}

Los requisitos funcionales describen las funcionalidades que el sistema debe ofrecer a los usuarios, junto con los criterios de aceptación que permiten verificar su correcta implementación.

A partir del análisis del problema y del backlog definido, se han identificado los siguientes requisitos funcionales principales:

\begin{itemize}
    \item \textbf{RF-01: Registro e inicio de sesión}: el sistema debe permitir a los usuarios registrarse, iniciar sesión y gestionar su cuenta.
    \textit{Criterio de aceptación}: un usuario puede crear una cuenta con correo y contraseña, iniciar sesión correctamente y cerrar sesión. Un intento con credenciales incorrectas muestra un mensaje de error.

    \item \textbf{RF-02: Gestión de perfil}: el usuario puede editar su información personal, incluyendo intereses, ubicación y cambio de contraseña, así como eliminar su cuenta.
    \textit{Criterio de aceptación}: los cambios en el perfil se reflejan inmediatamente en la aplicación. El usuario puede eliminar su cuenta y todos sus datos asociados son eliminados del sistema.

    \item \textbf{RF-03: Consulta de eventos}: el sistema debe mostrar un listado de eventos con información básica y permitir acceder al detalle de cada uno.
    \textit{Criterio de aceptación}: el listado muestra título, categoría, fecha y ubicación. Al pulsar sobre un evento se accede a su pantalla de detalle con toda la información disponible.

    \item \textbf{RF-04: Búsqueda y filtrado}: el usuario puede filtrar eventos por categoría, precio, distancia, fecha, título o descripción.
    \textit{Criterio de aceptación}: al aplicar uno o varios filtros, el listado se actualiza mostrando únicamente los eventos que cumplen los criterios seleccionados.

    \item \textbf{RF-05: Mapa de eventos}: el sistema debe mostrar los eventos en un mapa interactivo en función de la ubicación del usuario.
    \textit{Criterio de aceptación}: los eventos aparecen como marcadores en el mapa. Al pulsar sobre un marcador se muestra información básica del evento y se puede navegar a su detalle.

    \item \textbf{RF-06: Creación de eventos}: los usuarios pueden crear eventos, diferenciando entre eventos públicos y privados.
    \textit{Criterio de aceptación}: el formulario de creación permite seleccionar visibilidad. Un evento público queda en estado pendiente tras su creación; un evento privado es accesible inmediatamente por su creador.

    \item \textbf{RF-07: Moderación de eventos}: los eventos públicos deben ser aprobados por un moderador antes de ser visibles al crearse o editarse. El moderador puede editarlos antes de aprobarlos y también puede rechazarlos.
    \textit{Criterio de aceptación}: el moderador ve en su panel los eventos pendientes, puede aprobarlos, rechazarlos o modificarlos. Tras la resolución, el creador recibe una notificación con el resultado.

    \item \textbf{RF-08: Participación en eventos}: los usuarios pueden apuntarse a eventos públicos o acceder a eventos privados mediante enlace.
    \textit{Criterio de aceptación}: un usuario puede apuntarse a un evento activo y aparece en la lista de asistentes. Un evento privado solo es accesible si el usuario dispone del enlace de invitación.

    \item \textbf{RF-09: Chat y mensajería}: el sistema permite la comunicación entre usuarios mediante chat en eventos y mensajes privados directos.
    \textit{Criterio de aceptación}: los mensajes en el chat de un evento se reciben en tiempo real por todos los participantes. Los mensajes privados entre dos usuarios se entregan de forma inmediata y muestran indicador de leído.

    \item \textbf{RF-10: Gestión de imágenes}: los usuarios pueden subir y visualizar imágenes asociadas a eventos, chats y perfiles.
    \textit{Criterio de aceptación}: las imágenes se suben correctamente y se muestran en el contexto correspondiente. Se rechaza la subida de archivos que superen 5 MB o que no sean de tipo permitido.

    \item \textbf{RF-11: Valoraciones}: los usuarios pueden valorar eventos mediante puntuaciones y comentarios.
    \textit{Criterio de aceptación}: un usuario puede dejar una puntuación de 1 a 5 estrellas y un comentario opcional. La valoración media del evento se actualiza y es visible en el detalle del evento.

    \item \textbf{RF-12: Recomendaciones}: el sistema ofrece recomendaciones personalizadas en función de la actividad del usuario, su cercanía y fecha.
    \textit{Criterio de aceptación}: la pantalla de recomendaciones muestra eventos distintos a los ya vistos, ordenados según la puntuación calculada por el algoritmo en base a intereses, ubicación y actividad previa.

    \item \textbf{RF-13: Rutas de eventos}: el sistema recomienda rutas de eventos personalizadas en base a los gustos del usuario, la distancia y la fecha entre ellos. Además, los usuarios pueden crear sus propias rutas públicas y el resto puede valorarlas.
    \textit{Criterio de aceptación}: el sistema genera al menos una ruta recomendada cuando el usuario tiene ubicación e intereses definidos. Un usuario puede crear una ruta pública que es visible y valorable por otros usuarios.

    \item \textbf{RF-14: Notificaciones}: el sistema envía notificaciones sobre eventos cercanos, recordatorios de eventos a los que se apuntó el usuario y resultados de moderación al crear o editar un evento.
    \textit{Criterio de aceptación}: el usuario recibe notificaciones push cuando: (1) hay eventos cercanos a su ubicación, (2) se acerca la fecha de un evento al que está apuntado, o (3) un evento suyo es aprobado o rechazado por el moderador.

    \item \textbf{RF-15: Gestión de roles}: el sistema distingue entre usuario, administrador y moderador, ofreciendo funcionalidades específicas para cada uno.
    \textit{Criterio de aceptación}: un usuario sin permisos no puede acceder a las pantallas de administración ni moderación. Un administrador puede cambiar el rol de cualquier usuario.

    \item \textbf{RF-16: Gestión de categorías}: el administrador puede crear, editar y eliminar categorías de eventos.
    \textit{Criterio de aceptación}: las categorías creadas por el administrador aparecen disponibles inmediatamente en el formulario de creación de eventos y en los filtros de búsqueda.

    \item \textbf{RF-17: Gestión de usuarios}: el administrador puede editar el rol y eliminar usuarios.
    \textit{Criterio de aceptación}: el administrador puede cambiar el rol de un usuario desde el panel de administración. Un usuario eliminado no puede volver a iniciar sesión con esa cuenta.

    \item \textbf{RF-18: Organización de eventos}: los usuarios pueden guardar eventos, apuntarse a ellos y añadirlos al calendario para consultarlos posteriormente.
    \textit{Criterio de aceptación}: los eventos guardados aparecen en la sección correspondiente. Los eventos con fecha a los que el usuario está apuntado aparecen en el calendario en su fecha correcta.

    \item \textbf{RF-19: Actualización automática de eventos}: el sistema obtiene y actualiza eventos de fuentes externas de forma periódica mediante scraping.
    \textit{Criterio de aceptación}: los eventos importados aparecen en el listado como eventos aprobados. Si un evento ya existe, sus datos se actualizan en lugar de crear un duplicado. Los eventos sin ubicación válida no se almacenan.

    \item \textbf{RF-20: Edición de eventos}: los usuarios pueden editar los eventos que han creado.
    \textit{Criterio de aceptación}: el creador de un evento puede modificar sus datos. Si el evento es público, la edición vuelve a pasar por moderación antes de publicarse.

    \item \textbf{RF-21: Acceso a eventos privados}: los usuarios pueden acceder a los eventos privados mediante un enlace privado que comparte el creador.
    \textit{Criterio de aceptación}: un usuario que accede mediante el enlace correcto puede ver el detalle del evento privado. Un usuario sin el enlace no puede ver ni listar el evento.
\end{itemize}

\section{Requisitos no funcionales}

Los requisitos no funcionales describen las características de calidad que debe cumplir el sistema.

\begin{itemize}
    \item \textbf{RNF-01: Usabilidad}: la aplicación debe ofrecer una interfaz intuitiva y fácil de usar, adaptada a dispositivos móviles.
    
    \item \textbf{RNF-02: Rendimiento}: el sistema debe responder de forma rápida a las acciones del usuario, especialmente en la carga de eventos y mapas.
    
    \item \textbf{RNF-03: Escalabilidad}: la arquitectura debe permitir el crecimiento del sistema en número de usuarios y eventos.
    
    \item \textbf{RNF-04: Seguridad}: la autenticación debe ser segura, utilizando servicios como Firebase Auth, y protegiendo los datos del usuario.
    
    \item \textbf{RNF-05: Disponibilidad}: el sistema debe estar disponible de forma continua para garantizar el acceso a los usuarios.
    
    \item \textbf{RNF-06: Mantenibilidad}: el código debe estar organizado y documentado para facilitar futuras modificaciones.
    
    \item \textbf{RNF-07: Compatibilidad}: la aplicación debe funcionar tanto en dispositivos Android como iOS.
    
    \item \textbf{RNF-08: Integración}: el sistema debe integrarse correctamente con servicios externos como mapas, almacenamiento de imágenes y notificaciones.
\end{itemize}

\section{Reglas de negocio}

Las reglas de negocio definen las restricciones y condiciones del sistema, asegurando la coherencia de los datos y el correcto funcionamiento de las funcionalidades. Estas reglas de negocio se han definido evitando duplicar funcionalidades ya descritas en los requisitos funcionales, centrándose únicamente en restricciones, validaciones y condiciones del sistema.

\subsection{Gestión de usuarios y roles}

\begin{itemize}
    \item \textbf{RN-01}: Todo usuario debe estar autenticado para acceder a las funcionalidades principales del sistema.
    \item \textbf{RN-02}: El rol por defecto al registrarse es el de usuario.
    \item \textbf{RN-03}: Solo los administradores pueden modificar roles o eliminar usuarios.
\end{itemize}

\subsection{Gestión de eventos}

\begin{itemize}
    \item \textbf{RN-04}: Todo evento público debe pasar por un proceso de moderación antes de ser visible.
    \item \textbf{RN-05}: No se pueden aprobar o rechazar eventos que ya tengan ese mismo estado.
    \item \textbf{RN-06}: Un evento público solo será visible para los usuarios si su estado es aprobado.
    \item \textbf{RN-07}: Los eventos finalizados no permiten nuevas inscripciones.
    \item \textbf{RN-08}: Cuando un evento pasa de privado a público, debe someterse a moderación antes de ser visible.
\end{itemize}

\subsection{Eventos privados}

\begin{itemize}
    \item \textbf{RN-09}: Los eventos privados solo son accesibles mediante un enlace de invitación.
    \item \textbf{RN-10}: Solo el creador del evento privado puede compartir el enlace de acceso.
    \item \textbf{RN-11}: Los eventos privados no son visibles en el listado general de eventos.
    \item \textbf{RN-12}: Un evento privado solo puede ser consultado por su creador o usuarios invitados.
    \item \textbf{RN-13}: Los eventos privados no requieren moderación.

\end{itemize}

\subsection{Participación en eventos}

\begin{itemize}
    \item \textbf{RN-14}: Un usuario solo puede apuntarse a eventos activos.
    \item \textbf{RN-15}: Un usuario no puede apuntarse más de una vez al mismo evento.
    \item \textbf{RN-16}: Un usuario no puede guardar más de una vez el mismo evento.
    \item \textbf{RN-17}: Solo los eventos con fecha definida deben aparecer en el calendario del usuario.
\end{itemize}

\subsection{Comunicación y contenido}

\begin{itemize}
    \item \textbf{RN-18}: Las imágenes deben ser de contenido adecuado; el cumplimiento de esta política queda bajo la responsabilidad del usuario.
    \item \textbf{RN-19}: Las imágenes deben ser de tipo permitido (JPEG, PNG, WEBP, HEIC, HEIF).
    \item \textbf{RN-20}: Las imágenes no deben superar un tamaño máximo establecido de 5 MB.
\end{itemize}

\subsection{Valoraciones y reseñas}

\begin{itemize}
    \item \textbf{RN-21}: Las valoraciones deben estar comprendidas entre 1 y 5 estrellas.
\end{itemize}

\subsection{Recomendaciones y geolocalización}

\begin{itemize}
    \item \textbf{RN-22}: La ubicación del usuario debe estar definida para ofrecer funcionalidades basadas en proximidad.
\end{itemize}

\subsection{Scraping}

\begin{itemize}
    \item \textbf{RN-23}: Si un evento importado ya existe, el scraping debe actualizarlo en lugar de crear un duplicado.
    \item \textbf{RN-24}: Los eventos sin ubicación válida no deben almacenarse.
    \item \textbf{RN-25}: Los eventos importados mediante scraping se consideran públicos y se aprueban automáticamente.

\end{itemize}

\subsection{Restricciones de integridad de datos}

\begin{itemize}
    \item \textbf{RN-26}: Todo evento debe tener título, descripción y categoría.
    \item \textbf{RN-27}: Todo evento debe tener una ubicación geográfica válida.
    \item \textbf{RN-28}: La latitud y longitud deben ser valores numéricos y en el rango (latitud entre -90 y 90, longitud entre -180 y 180).
    \item \textbf{RN-29}: La fecha de inicio no puede ser posterior a la fecha de finalización.
    \item \textbf{RN-30}: El precio de un evento no puede ser negativo.
    \item \textbf{RN-31}: El radio de búsqueda debe ser un valor positivo.
    \item \textbf{RN-32}: Un evento no puede tener simultáneamente un precio fijo y un rango de precios mínimo y máximo.
    \item \textbf{RN-33}: La dirección de un evento debe ser coherente con sus coordenadas geográficas.
    \item \textbf{RN-34}: En un rango de precios, el precio mínimo no puede ser superior al precio máximo.


\end{itemize}

\subsection{Notificaciones}

\begin{itemize}
    \item \textbf{RN-35}: Las notificaciones de eventos cercanos solo pueden generarse si el usuario tiene una ubicación definida.
    \item \textbf{RN-36}: Las notificaciones de recordatorio solo pueden generarse para eventos a los que el usuario esté apuntado.
    \item \textbf{RN-37}: Las notificaciones de moderación solo deben enviarse al creador del evento aprobado o rechazado.
\end{itemize}

\section{Conclusiones}

El análisis del problema ha permitido identificar las necesidades principales que debe cubrir el sistema \textbf{Sevillaneando}, así como definir los requisitos necesarios para su desarrollo.

Se ha determinado que el sistema debe ir más allá de una simple plataforma de consulta de eventos, incorporando funcionalidades sociales, personalización y geolocalización para mejorar la experiencia del usuario.

La definición clara de las reglas de negocio, los requisitos de información, funcionales y no funcionales ha servido como base para el diseño e implementación del sistema, garantizando la coherencia entre el problema inicial y la solución desarrollada.